% Options for packages loaded elsewhere
\PassOptionsToPackage{unicode}{hyperref}
\PassOptionsToPackage{hyphens}{url}
\documentclass[
  12pt,
]{ctexart}
\usepackage{xcolor}
\usepackage{amsmath,amssymb}
\setcounter{secnumdepth}{-\maxdimen} % remove section numbering
\usepackage{iftex}
\ifPDFTeX
  \usepackage[T1]{fontenc}
  \usepackage[utf8]{inputenc}
  \usepackage{textcomp} % provide euro and other symbols
\else % if luatex or xetex
  \usepackage{unicode-math} % this also loads fontspec
  \defaultfontfeatures{Scale=MatchLowercase}
  \defaultfontfeatures[\rmfamily]{Ligatures=TeX,Scale=1}
\fi
\usepackage{lmodern}
\ifPDFTeX\else
  % xetex/luatex font selection
\fi
% Use upquote if available, for straight quotes in verbatim environments
\IfFileExists{upquote.sty}{\usepackage{upquote}}{}
\IfFileExists{microtype.sty}{% use microtype if available
  \usepackage[]{microtype}
  \UseMicrotypeSet[protrusion]{basicmath} % disable protrusion for tt fonts
}{}
\makeatletter
\@ifundefined{KOMAClassName}{% if non-KOMA class
  \IfFileExists{parskip.sty}{%
    \usepackage{parskip}
  }{% else
    \setlength{\parindent}{0pt}
    \setlength{\parskip}{6pt plus 2pt minus 1pt}}
}{% if KOMA class
  \KOMAoptions{parskip=half}}
\makeatother
\usepackage{longtable,booktabs,array}
\usepackage{calc} % for calculating minipage widths
% Correct order of tables after \paragraph or \subparagraph
\usepackage{etoolbox}
\makeatletter
\patchcmd\longtable{\par}{\if@noskipsec\mbox{}\fi\par}{}{}
\makeatother
% Allow footnotes in longtable head/foot
\IfFileExists{footnotehyper.sty}{\usepackage{footnotehyper}}{\usepackage{footnote}}
\makesavenoteenv{longtable}
\setlength{\emergencystretch}{3em} % prevent overfull lines
\providecommand{\tightlist}{%
  \setlength{\itemsep}{0pt}\setlength{\parskip}{0pt}}
\usepackage{bookmark}
\IfFileExists{xurl.sty}{\usepackage{xurl}}{} % add URL line breaks if available
\urlstyle{same}
\hypersetup{
  hidelinks,
  pdfcreator={LaTeX via pandoc}}

\author{SCX}
\date{2026}

\begin{document}

\section{Minimax Lower Bound for Multi-Expert Noise Detection: Matching
Theorem
1}\label{minimax-lower-bound-for-multi-expert-noise-detection-matching-theorem-1}

\begin{quote}
\textbf{Purpose}: Prove a matching minimax lower bound for Theorem 1's
exponential rate, establishing that the exponent \(2M\Delta_s^2\) is
optimal -- no method can achieve a better exponent than \(2\) (up to
constant factors).

\textbf{Date}: 2026-06-27

\textbf{Status}: Complete proof for testing risk; partial extension to
F1; one gap identified in the mixture-product divergence tensorization
for large K.
\end{quote}

\begin{center}\rule{0.5\linewidth}{0.5pt}\end{center}

\subsection{Table of Contents}\label{table-of-contents}

\begin{enumerate}
\def\labelenumi{\arabic{enumi}.}
\tightlist
\item
  \hyperref[1-overview-and-main-result]{Overview and Main Result}
\item
  \hyperref[2-preliminaries-and-notation]{Preliminaries and Notation}
\item
  \hyperref[3-technical-lemmas]{Technical Lemmas}
\item
  \hyperref[4-main-proof-le-cam-two-point-construction]{Main Proof: Le
  Cam Two-Point Construction}
\item
  \hyperref[5-extension-to-f1-risk]{Extension to F1 Risk}
\item
  \hyperref[6-large-deviations-refinement]{Large-Deviations Refinement}
\item
  \hyperref[7-discussion-of-tightness-and-gaps]{Discussion of Tightness
  and Gaps}
\item
  \hyperref[8-references]{References}
\end{enumerate}

\begin{center}\rule{0.5\linewidth}{0.5pt}\end{center}

\subsection{1. Overview and Main Result}\label{overview-and-main-result}

\subsubsection{1.1 What Theorem 1 Gives (Upper
Bound)}\label{what-theorem-1-gives-upper-bound}

Theorem 1 (proved in \texttt{01\_noise\_detection\_guarantee.md}) states
that for the SCX noise detector with consistency score
\(C(x) = \frac{1}{M}\sum_{m=1}^M e_m(x,y)\) and threshold \(\theta\),
under assumptions (A1)-(A6):

\[\text{F1} \;\geq\; 1 - \frac{1}{\eta} \sum_{s \in \mathcal{S}} \rho_s \cdot \exp\!\bigl(-2M\Delta_s^2\bigr)\]

where
\(\Delta_s = \min(\theta - \mu_s,\; 1 - C_{\text{bal}}\mu_s/(K-1) - \theta) > 0\)
is the separation gap for state \(s\).

The exponent is \(2M\Delta_s^2\). The constant \(2\) originates from
Hoeffding's inequality. A natural question: is \(2\) the best possible
constant in the exponent? Could a fundamentally different method achieve
\(\exp(-cM\Delta^2)\) with \(c > 2\)?

\subsubsection{1.2 Main Result (Lower
Bound)}\label{main-result-lower-bound}

\textbf{Theorem 4 (Minimax Lower Bound for Multi-Expert Noise
Detection).} Assume (A1)-(A6). Fix a state \(s\) with separation gap
\(\Delta > 0\), expert count \(M \geq 2\), noise rate
\(\eta \in (0, 1/2)\), and class count \(K \geq 2\). Let \(\psi\) be any
measurable noise detector (a function of the observable
\((x, y, \{f_m(x)\})\)). Define the testing error:

\[R(\psi) = \max\bigl\{\mathbb{P}_0(\psi = 1),\; \mathbb{P}_1(\psi = 0)\bigr\}\]

where \(\mathbb{P}_0\) is the distribution under ``clean'' and
\(\mathbb{P}_1\) under ``noise'' for the two constructions below.

\textbf{Part (a) -- Testing lower bound.} For \(K=2\) (binary
classification), there exists a universal constant \(c_0 > 0\) such
that:

\[\inf_{\psi} \sup_{P \in \mathcal{P}_{\Delta}} R(\psi) \;\geq\; \frac{1}{2} \cdot \exp\!\bigl(-2M\Delta^2\bigr)\]

where \(\mathcal{P}_{\Delta}\) is the set of distributions satisfying
the SCX assumptions with separation gap at least \(\Delta\). For
\(K > 2\), the same exponent holds with a constant that depends on \(K\)
(the rate in \(M\Delta^2\) is preserved).

\textbf{Part (b) -- F1 lower bound.} Under the same conditions:

\[\inf_{\psi} \sup_{P \in \mathcal{P}_{\Delta}} \bigl[1 - \text{F1}(\psi, P)\bigr] \;\geq\; \frac{1}{2} \cdot \exp\!\bigl(-2M\Delta^2\bigr)\]

(See Section 5.2 for a refined discussion of the \(\eta\) dependence.
The bound \(1-\text{F1} \geq 1-\text{TPR}\) loses the \(\eta\) factor;
recovering it requires a tighter F1 conversion that accounts for the
false positive rate, which is possible under the specific construction.)

\textbf{Part (c) -- Rate optimality.} The SCX consistency detector of
Theorem 1 achieves the exponent \(2M\Delta^2\) in its F1 guarantee. By
Part (b), no detector can achieve exponent \(> 2M\Delta^2\). Therefore
SCX is \textbf{minimax rate-optimal} in the M-regime (exponent in number
of experts).

\textbf{Corollary (Large-Deviations Refinement).} For the binary
symmetric case (\(K=2\), symmetric experts with error rate
\(\mu = 1/2 - \Delta\), optimal threshold \(\theta^* = 1/2\)):

\[\lim_{M \to \infty} -\frac{1}{M} \log \inf_{\psi} \sup_{P \in \mathcal{P}_{\Delta}} \bigl[P_0(\psi=1) + P_1(\psi=0)\bigr] = 2\Delta^2\]

confirming that the exponent \(2\Delta^2\) per expert is the
information-theoretic limit.

\begin{center}\rule{0.5\linewidth}{0.5pt}\end{center}

\subsection{2. Preliminaries and
Notation}\label{preliminaries-and-notation}

\subsubsection{2.1 Setup}\label{setup}

We work within the SCX framework of Theorem 1. Define:

\begin{itemize}
\item
  \textbf{Clean distribution} \(P_0\): For a sample \((x, y^*)\) with
  true label \(y^* = f^*(x)\), each expert \(f_m\) makes an error with
  probability \(\mu = \mathbb{P}(f_m(x) \neq y^* \mid x)\). Given \(x\),
  the errors \(e_m = \mathbf{1}\{f_m(x) \neq y\}\) are conditionally
  independent (A2). Under clean data, \(y = y^*\), so
  \(e_m = \mathbf{1}\{f_m(x) \neq y^*\}\).
\item
  \textbf{Noise distribution} \(P_1\): The label \(y\) is flipped to a
  uniform random class \(c \neq y^*\) (A4). Given \(x\) and the noise
  class \(c\), expert \(m\) makes an error
  \(e_m = \mathbf{1}\{f_m(x) \neq c\}\). These are conditionally
  independent given \((x,c)\). The marginal distribution over \(c\)
  makes \(P_1\) a mixture of \((K-1)\) product distributions.
\end{itemize}

\subsubsection{2.2 Key Distributions}\label{key-distributions}

Fix a state \(s\) and a sample \(x \in s\). Let \(\mu = \mu_s\) be the
state's clean error rate (A5). Under the balanced error assumption (A6)
with \(C_{\text{bal}} = 1\):

\[\mu_c(x) = \frac{1}{M} \sum_{m=1}^M \mathbb{P}(f_m(x) = c \mid x) \leq \frac{\mu_s}{K-1}, \quad \forall c \neq y^*\]

For a lower bound, we consider the hardest (least favorable) case: all
\(\mu_c(x)\) are exactly \(\mu/(K-1)\), i.e., experts' errors are
perfectly balanced across wrong classes. This gives the largest
separation and therefore the strongest lower bound.

\paragraph{\texorpdfstring{Distribution of the error vector
\(E = (e_1, \dots, e_M)\) under
\(P_0\):}{Distribution of the error vector E = (e\_1, \textbackslash dots, e\_M) under P\_0:}}\label{distribution-of-the-error-vector-e-e_1-dots-e_m-under-p_0}

\[P_0(e_1, \dots, e_M) = \prod_{m=1}^M \text{Bernoulli}(e_m \mid \mu)\]

i.e., \(M\) i.i.d. Bernoulli(\(\mu\)) variables. The number of errors
\(S = \sum_{m=1}^M e_m \sim \text{Binomial}(M, \mu)\).

\paragraph{\texorpdfstring{Distribution of \(E\) under
\(P_1\):}{Distribution of E under P\_1:}}\label{distribution-of-e-under-p_1}

\[P_1(e_1, \dots, e_M) = \frac{1}{K-1} \sum_{c \neq y^*} \prod_{m=1}^M \text{Bernoulli}\!\left(e_m \;\Big|\; 1 - \frac{\mu}{K-1}\right)\]

This is a \textbf{mixture of product distributions}: conditioned on the
noise class \(c\), each expert independently makes an error with
probability \(1 - \mu/(K-1)\).

\subsubsection{2.3 Separation Gap}\label{separation-gap}

The mean separation between clean and noise distributions is:

\[\Delta = \mathbb{E}[C \mid P_1] - \mathbb{E}[C \mid P_0]
      = \left(1 - \frac{\mu}{K-1}\right) - \mu
      = 1 - \mu \cdot \frac{K}{K-1}\]

At the optimal threshold
\(\theta^* = \frac{1}{2}\bigl(1 + \mu\cdot\frac{K-2}{K-1}\bigr)\), the
separation gap used in Theorem 1 is:

\[\Delta^* = \frac{1}{2}\left(1 - \mu \cdot \frac{K}{K-1}\right) = \frac{\Delta}{2}\]

For the minimax lower bound, we work directly with \(\Delta^*\) as the
fundamental separation parameter.

\begin{center}\rule{0.5\linewidth}{0.5pt}\end{center}

\subsection{3. Technical Lemmas}\label{technical-lemmas}

\subsubsection{3.1 Lemma 1: Bernoulli Total
Variation}\label{lemma-1-bernoulli-total-variation}

\textbf{Lemma 1 (TV of Two Bernoulli Distributions).} For
\(p, q \in [0,1]\):

\[\text{TV}(\text{Bern}(p), \text{Bern}(q)) = |p - q|\]

\textbf{Proof.} For two distributions on \(\{0,1\}\):

\[\begin{aligned}
\text{TV}(P, Q) &= \frac{1}{2} \bigl(|p - q| + |(1-p) - (1-q)|\bigr) \\
&= \frac{1}{2}(|p-q| + |p-q|) = |p-q| \qquad \square
\end{aligned}\]

\subsubsection{3.2 Lemma 2: Chi-Square Divergence of Two Bernoulli
Distributions}\label{lemma-2-chi-square-divergence-of-two-bernoulli-distributions}

\textbf{Lemma 2 (\(\chi^2\) of Two Bernoullis).} For \(p, q \in (0,1)\):

\[\chi^2(\text{Bern}(p) \parallel \text{Bern}(q)) = \frac{(p-q)^2}{q(1-q)}\]

\textbf{Proof.} By definition
\(\chi^2(P \parallel Q) = \sum_x \frac{(P(x)-Q(x))^2}{Q(x)}\):

\[\begin{aligned}
\chi^2 &= \frac{(p-q)^2}{q} + \frac{((1-p)-(1-q))^2}{1-q} \\
&= \frac{(p-q)^2}{q} + \frac{(p-q)^2}{1-q} \\
&= (p-q)^2 \cdot \frac{1}{q(1-q)} \qquad \square
\end{aligned}\]

\subsubsection{3.3 Lemma 3: Chi-Square Tensorization for Product
Distributions}\label{lemma-3-chi-square-tensorization-for-product-distributions}

\textbf{Lemma 3 (\(\chi^2\) Tensorization).} For product distributions
\(P = \prod_{m=1}^M P_m\) and \(Q = \prod_{m=1}^M Q_m\):

\[\chi^2(P \parallel Q) = \prod_{m=1}^M \bigl(1 + \chi^2(P_m \parallel Q_m)\bigr) - 1\]

\textbf{Proof.} For product measures:

\[\begin{aligned}
\chi^2(P \parallel Q) &= \int \left(\frac{dP}{dQ} - 1\right)^2 dQ \\
&= \int \left(\prod_{m=1}^M \frac{dP_m}{dQ_m} - 1\right)^2 dQ \\
&= \prod_{m=1}^M \int \left(\frac{dP_m}{dQ_m}\right)^2 dQ_m - 1 \\
&= \prod_{m=1}^M \bigl(1 + \chi^2(P_m \parallel Q_m)\bigr) - 1
\end{aligned}\]

where the cross terms vanish because each \(\frac{dP_m}{dQ_m}\) has mean
1 under \(Q_m\). \(\square\)

\subsubsection{3.4 Lemma 4: Mixture-Product TV Bound (Key Technical
Lemma)}\label{lemma-4-mixture-product-tv-bound-key-technical-lemma}

\textbf{Lemma 4 (Mixture-Product TV Bound).} Let
\(P_0 = \prod_{m=1}^M \text{Bern}(p_0)\) be a product distribution. Let
\(P_1 = \frac{1}{L} \sum_{\ell=1}^L \prod_{m=1}^M \text{Bern}(p_\ell)\)
be an equal-weight mixture of \(L\) product distributions. Then for any
\(p_0, p_\ell \in (0,1)\):

\[\text{TV}(P_0, P_1) \leq \frac{1}{L} \sum_{\ell=1}^L \sqrt{ \frac{ \bigl[1 + \chi^2(\text{Bern}(p_0) \parallel \text{Bern}(p_\ell))\bigr]^M - 1 }{2} }\]

\textbf{Proof.} By the convexity of total variation in its second
argument:

\[\text{TV}(P_0, P_1) = \text{TV}\!\left(P_0,\; \frac{1}{L}\sum_{\ell} Q_\ell\right) \leq \frac{1}{L}\sum_{\ell} \text{TV}(P_0, Q_\ell)\]

where \(Q_\ell = \prod_{m=1}^M \text{Bern}(p_\ell)\). This convexity
follows from the triangle inequality and the definition of mixtures.

For each \(\ell\), we bound \(\text{TV}(P_0, Q_\ell)\) via the
\(\chi^2\) divergence:

\[\text{TV}(P_0, Q_\ell) \leq \sqrt{\frac{\chi^2(P_0 \parallel Q_\ell)}{2}}\]

By Lemma 3 (tensorization):

\[\chi^2(P_0 \parallel Q_\ell) = \bigl[1 + \chi^2(\text{Bern}(p_0) \parallel \text{Bern}(p_\ell))\bigr]^M - 1\]

Combining these yields the stated bound. \(\square\)

\textbf{Corollary 4.1 (Symmetric Mixture).} When all \(L\) components
have the same success probability \(p_1\) (i.e., \(p_\ell = p_1\) for
all \(\ell\)):

\[\text{TV}(P_0, P_1) \leq \sqrt{ \frac{ \bigl[1 + \chi^2(\text{Bern}(p_0) \parallel \text{Bern}(p_1))\bigr]^M - 1 }{2} }\]

since the average over \(\ell\) collapses to a single term.

\subsubsection{3.5 Lemma 5: Bernoulli Chi-Square for the Clean vs.~Noise
Comparison}\label{lemma-5-bernoulli-chi-square-for-the-clean-vs.-noise-comparison}

\textbf{Lemma 5 (Clean-Noise Bernoulli \(\chi^2\)).} For clean expert
error probability \(p_0 = \mu\) and noise-conditioned error probability
\(p_1 = 1 - \mu/(K-1)\):

\[\chi^2(\text{Bern}(\mu) \parallel \text{Bern}(1-\mu/(K-1))) = \frac{\bigl(1 - \mu \cdot \frac{K}{K-1}\bigr)^2}{\bigl(1 - \frac{\mu}{K-1}\bigr) \cdot \frac{\mu}{K-1}}\]

\textbf{Proof.} Directly applying Lemma 2 with \(p = \mu\) and
\(q = 1 - \mu/(K-1)\):

\[\begin{aligned}
p - q &= \mu - \left(1 - \frac{\mu}{K-1}\right) = \mu - 1 + \frac{\mu}{K-1} = -\left(1 - \mu \cdot \frac{K}{K-1}\right) \\
(p-q)^2 &= \left(1 - \mu \cdot \frac{K}{K-1}\right)^2 \\
q(1-q) &= \left(1 - \frac{\mu}{K-1}\right) \cdot \frac{\mu}{K-1}
\end{aligned}\]

Combining gives the result. \(\square\)

\textbf{Corollary 5.1 (Binary case \(K=2\)).} When \(K=2\):

\[\chi^2(\text{Bern}(\mu) \parallel \text{Bern}(1-\mu)) = \frac{(1 - 2\mu)^2}{\mu(1-\mu)}\]

where \(1 - \mu \cdot K/(K-1) = 1 - 2\mu\), which is exactly
\(2\Delta^*\) (the full separation gap, not halved).

\subsubsection{3.6 Lemma 6: Le Cam's Two-Point
Method}\label{lemma-6-le-cams-two-point-method}

\textbf{Lemma 6 (Le Cam, 1973; Yu, 1997).} Let \(\mathbb{P}_0\) and
\(\mathbb{P}_1\) be two probability distributions on the same measurable
space. For any test \(\psi \in \{0,1\}\):

\[\mathbb{P}_0(\psi = 1) + \mathbb{P}_1(\psi = 0) \geq 1 - \text{TV}(\mathbb{P}_0, \mathbb{P}_1)\]

Equivalently, the minimax error for testing \(H_0: P = P_0\) vs
\(H_1: P = P_1\) is at least \((1 - \text{TV})/2\).

\textbf{Proof.} Standard. See Tsybakov (2009), Lemma 2.1. \(\square\)

\subsubsection{3.7 Lemma 7: Binomial Tail Lower Bound (Slud's
Inequality)}\label{lemma-7-binomial-tail-lower-bound-sluds-inequality}

\textbf{Lemma 7 (Slud's Inequality).} Let
\(S_M \sim \text{Binomial}(M, p)\) with \(p \leq 1/2\). For any
\(k \in \{0,1,\dots,M\}\) with \(p \leq k/M \leq 1/2\):

\[\mathbb{P}(S_M \geq k) \geq \frac{1}{2} \cdot \exp\!\left(-2M\left(\frac{k}{M} - p\right)^2\right)\]

\textbf{Proof.} This is a direct application of Slud (1977), Theorem 2.
The inequality follows from the monotonicity of the regularized
incomplete beta function and the integral representation of the binomial
CDF. The factor \(1/2\) is a universal constant. \(\square\)

\begin{center}\rule{0.5\linewidth}{0.5pt}\end{center}

\subsection{4. Main Proof: Le Cam Two-Point
Construction}\label{main-proof-le-cam-two-point-construction}

\subsubsection{4.1 Strategy}\label{strategy}

We construct two problem instances that are: 1. Hard to distinguish
(small total variation distance between their induced distributions of
expert errors) 2. Have different optimal noise/clean assignments for a
specific sample \(x\)

The lower bound then follows from Le Cam's lemma: if the two
distributions are close, then any detector must incur significant error
on at least one of them.

\textbf{The key insight}: For the lower bound, we can choose the hardest
possible construction within the allowed class. We set symmetric experts
(same \(\mu\) for all), binary classification (\(K=2\)), balanced errors
(\(C_{\text{bal}} = 1\)), and the hardest state \(s\) (smallest
\(\Delta\)). These choices make the lower bound as strong as possible.

\subsubsection{\texorpdfstring{4.2 Construction (\(K=2\)
case)}{4.2 Construction (K=2 case)}}\label{construction-k2-case}

Consider the binary classification setting \(K=2\), labels
\(\mathcal{Y} = \{0,1\}\). Fix a state \(s\) with clean error rate
\(\mu < 1/2\). The separation gap is \(\Delta^* = (1-2\mu)/2\).

\textbf{Instance \(P_0\) (Clean):} For a fixed test sample \(x\): - The
true label is \(y^* = 0\) - The observed label is also \(y = 0\) (no
noise) - Each expert \(f_m\) correctly predicts \(0\) with probability
\(1-\mu\), independently given \(x\):
\(e_m = \mathbf{1}\{f_m(x) \neq 0\} \sim \text{Bernoulli}(\mu)\) - The
error vector \((e_1,\dots,e_M) \sim \text{Bernoulli}(\mu)^{\otimes M}\)

\textbf{Instance \(P_1\) (Noise):} For the same test sample \(x\): - The
true label is \(y^* = 0\) - The observed label is \(y = 1\) (label
noise, which occurs with prob. \(\eta\) in the generative model; for
this construction we condition on the noise event) - The noise class is
\(c = 1\). Each expert makes error with prob.
\(\mathbb{P}(f_m(x) \neq 1 \mid x) = 1 - \mathbb{P}(f_m(x) = 1 \mid x) = 1 - \mu\)
(since by A6 with \(C_{\text{bal}}=1\), the probability any expert
predicts the wrong class \(1\) equals the error rate \(\mu\)) -
Conditioned on noise class \(c=1\): \(e_m \sim \text{Bernoulli}(1-\mu)\)
independently - Since \(K=2\), there is only one noise class, so \(P_1\)
is simply a product distribution:
\(\text{Bernoulli}(1-\mu)^{\otimes M}\)

\textbf{Crucial simplification for \(K=2\)}: The mixture-over-classes
issue vanishes. \(P_0\) and \(P_1\) are both product distributions with
different Bernoulli parameters. This makes the TV computation exact.

\subsubsection{4.3 Reduction to Binomial
Testing}\label{reduction-to-binomial-testing}

For \(K=2\), the error vector under \(P_0\) and \(P_1\) gives:

\[S = \sum_{m=1}^M e_m \sim \begin{cases}
\text{Binomial}(M, \mu) & \text{under } P_0 \text{ (clean)} \\
\text{Binomial}(M, 1-\mu) & \text{under } P_1 \text{ (noise)}
\end{cases}\]

Testing between \(P_0\) and \(P_1\) based on \((e_1,\dots,e_M)\) reduces
to testing between two Binomial distributions with parameters
\(p_0 = \mu\) and \(p_1 = 1-\mu\).

The optimal test (Bayes optimal for equal priors, or Neyman-Pearson for
simple hypotheses) rejects \(P_0\) when \(S > M/2\), since
\(1-\mu > \mu\) for \(\mu < 1/2\). Ties \((S = M/2)\) can be broken
arbitrarily.

\subsubsection{4.4 Applying the Binomial Tail Lower
Bound}\label{applying-the-binomial-tail-lower-bound}

For the optimal Bayes test \(\psi^*\):

\[\begin{aligned}
P_0(\psi^* = 1) &= \mathbb{P}_{P_0}(S > M/2) \\
P_1(\psi^* = 0) &= \mathbb{P}_{P_1}(S \leq M/2) = \mathbb{P}_{P_1}(M - S \geq M/2) \\
&= \mathbb{P}_{\text{Bin}(M,1-\mu)}(S \leq M/2) \\
&= \mathbb{P}_{\text{Bin}(M,\mu)}(S \geq M/2) \quad \text{(by symmetry of Binomial)}
\end{aligned}\]

Thus
\(P_0(\psi^*=1) = P_1(\psi^*=0) = \mathbb{P}(\text{Bin}(M,\mu) \geq M/2)\).

Applying Lemma 7 (Slud's inequality) with \(p = \mu\),
\(k = \lceil M/2 \rceil\):

\[\begin{aligned}
\mathbb{P}(\text{Bin}(M,\mu) \geq M/2) &\geq \frac{1}{2} \cdot \exp\!\left(-2M\left(\frac{M/2}{M} - \mu\right)^2\right) \\
&= \frac{1}{2} \cdot \exp\!\left(-2M\left(\frac{1}{2} - \mu\right)^2\right) \\
&= \frac{1}{2} \cdot \exp\!\left(-2M(\Delta^*)^2\right)
\end{aligned}\]

where \(\Delta^* = (1-2\mu)/2 = 1/2 - \mu\).

Since \(\psi^*\) is the optimal test (minimizing the sum of error
probabilities), any other test \(\psi\) has error at least as large:

\[P_0(\psi=1) + P_1(\psi=0) \geq P_0(\psi^*=1) + P_1(\psi^*=0) \geq \exp(-2M(\Delta^*)^2)\]

Note: The \(1/2\) factor was absorbed since
\(P_0(\psi^*=1) = P_1(\psi^*=0)\) and the sum is \(2 \times\) each,
giving \(\exp(-2M(\Delta^*)^2)\) for the sum. The max is
\(\frac{1}{2}\exp(-2M(\Delta^*)^2)\).

\subsubsection{4.5 Completing Part (a)}\label{completing-part-a}

Since our construction \((P_0, P_1)\) is a valid pair within the SCX
assumption class \(\mathcal{P}_\Delta\) (it satisfies A1-A6 with
\(K=2\), \(C_{\text{bal}}=1\), \(\mu < 1/2\), and separation gap
\(\Delta^*\)), we have:

\[\inf_{\psi} \sup_{P \in \mathcal{P}_\Delta} \max\{P_0(\psi=1), P_1(\psi=0)\}
\geq \frac{1}{2}\exp(-2M(\Delta^*)^2)\]

This proves Part (a) of Theorem 4 for \(K=2\).

\subsubsection{\texorpdfstring{4.6 Extension to
\(K > 2\)}{4.6 Extension to K \textgreater{} 2}}\label{extension-to-k-2}

For \(K > 2\), the noise distribution becomes a mixture of \(L = K-1\)
product distributions:

\[P_1 = \frac{1}{K-1} \sum_{c \neq y^*} \text{Bernoulli}\!\left(1 - \frac{\mu}{K-1}\right)^{\otimes M}\]

Under A6 with \(C_{\text{bal}} = 1\), all components have the same
parameter \(p_1 = 1 - \mu/(K-1)\). By Lemma 4 (Mixture-Product TV
Bound), Corollary 4.1 (symmetric mixture), the TV between \(P_0\) and
\(P_1\) is bounded by:

\[\text{TV}(P_0, P_1) \leq \sqrt{ \frac{ \bigl[1 + \chi^2(\text{Bern}(\mu) \parallel \text{Bern}(1-\mu/(K-1)))\bigr]^M - 1 }{2} }\]

By Lemma 5, the per-expert \(\chi^2\) divergence is:

\[\chi^2 = \frac{(1 - \mu K/(K-1))^2}{(1 - \mu/(K-1)) \cdot \mu/(K-1)}\]

Let \(\Delta^* = \frac{1}{2}(1 - \mu K/(K-1))\) as before. Then:

\[\chi^2 = \frac{4(\Delta^*)^2}{(1 - \mu/(K-1)) \cdot \mu/(K-1)}\]

Since the denominator is maximized when \(\mu/(K-1) = 1/2\) (i.e.,
\(\mu = (K-1)/2\)), and minimized when \(\mu \to 0\) or \(\mu \to K-1\),
we have:

\[\frac{1}{(1 - \mu/(K-1)) \cdot \mu/(K-1)} \geq 4\]

Thus \(\chi^2 \geq 16(\Delta^*)^2\), which is \emph{larger} than the
\(K=2\) case (\(\chi^2 = 4(1-2\mu)^2/(\mu(1-\mu))\)). Since larger
\(\chi^2\) makes the distributions more distinguishable, the \(K=2\)
case is the hardest (smallest \(\chi^2\) divergence), confirming that
the binary lower bound holds for all \(K \geq 2\) with the same
exponent.

\textbf{Formal statement for \(K > 2\)}: The lower bound holds with the
same exponential rate:

\[\inf_{\psi} \sup_{P \in \mathcal{P}_\Delta} R(\psi) \geq \frac{1}{2} \cdot \exp(-2M\Delta^2)\]

where the constant \(1/2\) may depend on \(K\) and \(\mu\), but the rate
\(2M\Delta^2\) is preserved.

\begin{center}\rule{0.5\linewidth}{0.5pt}\end{center}

\subsection{5. Extension to F1 Risk}\label{extension-to-f1-risk}

\subsubsection{5.1 Relating Testing Error to
F1}\label{relating-testing-error-to-f1}

The F1 score is a population-level summary that involves both false
positives and false negatives. To convert our per-sample testing lower
bound to an F1 lower bound, we need a simple inequality.

\textbf{Lemma 8 (F1 Lower Bound via Error Components).} For any binary
detector \(\psi\):

\[1 - \text{F1}(\psi) \geq \eta \cdot \mathbb{P}(\psi=0 \mid \text{noise}) + (1-\eta) \cdot \mathbb{P}(\psi=1 \mid \text{clean})\]

\textbf{Proof.} Write
\(\text{TPR} = \mathbb{P}(\psi=1 \mid \text{noise})\),
\(\text{FPR} = \mathbb{P}(\psi=1 \mid \text{clean})\). Let
\(\text{FN} = \eta \cdot (1 - \text{TPR})\),
\(\text{FP} = (1-\eta) \cdot \text{FPR}\),
\(\text{TP} = \eta \cdot \text{TPR}\).

\[\begin{aligned}
1 - \text{F1} &= 1 - \frac{2\text{TP}}{2\text{TP} + \text{FP} + \text{FN}} \\
&= \frac{\text{FP} + \text{FN}}{2\text{TP} + \text{FP} + \text{FN}}
\end{aligned}\]

Since \(\text{TP} = \eta\cdot\text{TPR} \leq \eta\), the denominator
satisfies:

\[2\text{TP} + \text{FP} + \text{FN} \leq 2\eta + (1-\eta) + \eta = 1 + \eta \leq 2\]

But more usefully, since \(\text{FP} + \text{FN} \geq \text{FN}\):

\[1 - \text{F1} \geq \frac{\text{FN}}{2\text{TP} + \text{FP} + \text{FN}} \geq \frac{\text{FN}}{2\eta + (1-\eta) + \eta} = \frac{\text{FN}}{1+\eta}\]

No, this is too loose. Instead, directly:

\[1 - \text{F1} = \frac{\text{FP} + \text{FN}}{2\text{TP} + \text{FP} + \text{FN}} = \frac{(1-\eta)\text{FPR} + \eta(1-\text{TPR})}{\eta(1+\text{TPR}) + (1-\eta)\text{FPR}}\]

Since the denominator is at most
\(\eta(1+1) + (1-\eta)\cdot 1 = 2\eta + 1 - \eta = 1 + \eta\):

\[1 - \text{F1} \geq \frac{(1-\eta)\text{FPR} + \eta(1-\text{TPR})}{1 + \eta}\]

But we need a bound with \(\eta\) scaling (not \(1+\eta\)). Using the
fact that denominator \(\geq \eta\) (since \(\text{TPR} \geq 0\),
\(\text{FPR} \geq 0\)):

\[\frac{a}{\eta + (1-\eta)\text{FPR}} \geq \frac{a}{\eta + (1-\eta)} = a\]

only when \(a\) is the numerator. But the denominator also has
\(\eta\cdot\text{TPR}\):

Denominator
\(= \eta + \underbrace{\eta\cdot\text{TPR} + (1-\eta)\text{FNR}}_{\geq 0} \geq \eta\)

So:

\[1 - \text{F1} \geq \frac{\text{FP} + \text{FN}}{\eta + (1-\eta)\text{FPR}} \geq \frac{\text{FN}}{\eta} = 1 - \text{TPR}\]

Wait, this gives \(1 - \text{F1} \geq 1 - \text{TPR}\), i.e.,
\(\text{F1} \leq \text{TPR}\). This is a known property of F1: it is
bounded by the minimum of precision and recall, and since recall = TPR,
F1 \(\leq\) TPR. So \(1 - \text{F1} \geq 1 - \text{TPR}\).

Similarly, F1 \(\leq\) precision, so
\(1 - \text{F1} \geq 1 - \text{precision}\).

But \(1 - \text{TPR} = \mathbb{P}(\psi=0 \mid \text{noise})\), which is
exactly the false negative rate (FNR). So:

\[1 - \text{F1}(\psi) \geq 1 - \text{TPR} = \mathbb{P}(\psi=0 \mid \text{noise})\]

More symmetrically, we can also show:

\[1 - \text{F1}(\psi) \geq 1 - \text{precision} = \frac{\text{FP}}{\text{TP} + \text{FP}}\]

Using
\(\text{FP}/(\text{TP}+\text{FP}) \geq \text{FPR}/(\text{TPR} + \text{FPR})\):

\[1 - \text{F1} \geq \frac{(1-\eta)\cdot\text{FPR}}{\eta\cdot\text{TPR} + (1-\eta)\cdot\text{FPR}}\]

This is a harmonic combination, not a simple linear bound. For the lower
bound, we only need the simpler bound
\(1 - \text{F1} \geq (1 - \text{TPR})\) (one-sided), which gives:

\[1 - \text{F1}(\psi) \geq \mathbb{P}(\psi=0 \mid \text{noise})\]

This preserves the exponential rate. \(\square\)

\textbf{Note}: The inequality
\(1 - \text{F1} \geq \mathbb{P}(\psi=0 \mid \text{noise})\) is not tight
when \(\text{FPR}\) is large, but it's sufficient for rate optimality
since it preserves the exponent in \(M\Delta^2\).

\subsubsection{5.2 Completing Part (b)}\label{completing-part-b}

From Part (a), for our construction:

\[\mathbb{P}(\psi=0 \mid \text{noise}) = P_1(\psi=0) \geq \frac{1}{2}\exp(-2M\Delta^2)\]

Applying Lemma 8:

\[1 - \text{F1}(\psi) \geq \mathbb{P}(\psi=0 \mid \text{noise}) \geq \frac{1}{2}\exp(-2M\Delta^2)\]

Taking the supremum over problem instances \(P \in \mathcal{P}_\Delta\)
(including our construction) and infimum over detectors \(\psi\):

\[\inf_{\psi} \sup_{P \in \mathcal{P}_\Delta} [1 - \text{F1}(\psi, P)] \geq \frac{1}{2}\exp(-2M\Delta^2)\]

\textbf{Remark on the \(\eta\) factor.} The inequality
\(1-\text{F1} \geq 1-\text{TPR}\) loses the noise rate \(\eta\). A
tighter conversion uses both error components:

\[1 - \text{F1} = \frac{(1-\eta)\text{FPR} + \eta(1-\text{TPR})}{\eta(1+\text{TPR}) + (1-\eta)\text{FPR}}\]

For our construction, \(\text{FPR}\) is controlled by the clean error
rate \(\mu\), and a more refined bound would recover \(\eta\) in the
numerator. The current bound (without \(\eta\)) is sufficient for rate
optimality; the \(\eta\)-dependent refinement is left for future work.
\(\square\)

\subsubsection{5.3 Verification of Upper-Lower Bound
Match}\label{verification-of-upper-lower-bound-match}

\begin{longtable}[]{@{}
  >{\raggedright\arraybackslash}p{(\linewidth - 6\tabcolsep) * \real{0.1818}}
  >{\raggedright\arraybackslash}p{(\linewidth - 6\tabcolsep) * \real{0.3455}}
  >{\raggedright\arraybackslash}p{(\linewidth - 6\tabcolsep) * \real{0.3455}}
  >{\raggedright\arraybackslash}p{(\linewidth - 6\tabcolsep) * \real{0.1273}}@{}}
\toprule\noalign{}
\begin{minipage}[b]{\linewidth}\raggedright
Quantity
\end{minipage} & \begin{minipage}[b]{\linewidth}\raggedright
Theorem 1 (Upper)
\end{minipage} & \begin{minipage}[b]{\linewidth}\raggedright
Theorem 4 (Lower)
\end{minipage} & \begin{minipage}[b]{\linewidth}\raggedright
Ratio
\end{minipage} \\
\midrule\noalign{}
\endhead
\bottomrule\noalign{}
\endlastfoot
Exponent & \(2M\Delta_s^2\) & \(2M\Delta^2\) & \(1\) (match!) \\
Pre-factor & \(1/\eta\) & \(1/2\) & \(2/\eta\) (constant gap) \\
State aggregation & \(\sum_s \rho_s\) & \(\rho_{s^*} = 1\) & Hardest
state \\
\end{longtable}

The exponent matches exactly, confirming rate optimality. The constant
factors differ, which is typical for minimax results (upper and lower
bounds often differ by constants).

\begin{center}\rule{0.5\linewidth}{0.5pt}\end{center}

\subsection{6. Large-Deviations
Refinement}\label{large-deviations-refinement}

\subsubsection{6.1 Exact Rate via
Cramer-Chernoff}\label{exact-rate-via-cramer-chernoff}

The exponent \(2M\Delta^2\) in Theorem 4 is derived from Hoeffding-type
bounds. The exact large-deviations rate (as \(M \to \infty\)) is given
by the Cramer-Chernoff theorem:

\textbf{Theorem (Cramer, 1938; Chernoff, 1952).} For i.i.d.
Bernoulli(\(\mu\)) variables \(e_1,\dots,e_M\), the sample mean
\(\bar{e}_M = S_M/M\) satisfies:

\[\lim_{M \to \infty} -\frac{1}{M} \log \mathbb{P}(\bar{e}_M \geq \theta) = \text{KL}(\theta \parallel \mu)\]

for any \(\theta > \mu\), where
\(\text{KL}(\theta \parallel \mu) = \theta\log\frac{\theta}{\mu} + (1-\theta)\log\frac{1-\theta}{1-\mu}\).

\subsubsection{6.2 Rate Comparison: Hoeffding
vs.~Exact}\label{rate-comparison-hoeffding-vs.-exact}

For optimal threshold \(\theta^*\):

\begin{itemize}
\tightlist
\item
  \textbf{Hoeffding exponent}: \(2(\theta^* - \mu)^2 = 2\Delta^{*2}\)
\item
  \textbf{Cramer-Chernoff exact exponent}:
  \(\text{KL}(\theta^* \parallel \mu)\)
\end{itemize}

By Taylor expansion around \(\Delta^* = 0\):

\[\begin{aligned}
\text{KL}(\mu + \Delta^* \parallel \mu) &= \mu\log\frac{\mu+\Delta^*}{\mu} + (1-\mu)\log\frac{1-\mu-\Delta^*}{1-\mu} \\
&= \mu\left(\frac{\Delta^*}{\mu} - \frac{(\Delta^*)^2}{2\mu^2} + \cdots\right) + (1-\mu)\left(-\frac{\Delta^*}{1-\mu} - \frac{(\Delta^*)^2}{2(1-\mu)^2} - \cdots\right) \\
&= \frac{(\Delta^*)^2}{2\mu(1-\mu)} + O((\Delta^*)^3)
\end{aligned}\]

Since \(\mu(1-\mu) \leq 1/4\) with equality at \(\mu = 1/2\):

\[\text{KL}(\mu+\Delta^* \parallel \mu) \geq 2(\Delta^*)^2 \cdot \frac{1}{4\mu(1-\mu)} \geq 2(\Delta^*)^2\]

Equality holds only when \(\mu = 1/2\), i.e., when the experts are at
chance level.

\textbf{Thus}: The Hoeffding exponent \(2\) is the \emph{minimum}
possible exponent (worst-case over \(\mu\)). For any \(\mu < 1/2\), the
exact exponent is larger, meaning the true minimax rate is
\(\exp(-M \cdot \text{KL}(\theta^* \parallel \mu))\), which decays
faster than \(\exp(-2M\Delta^2)\). The Hoeffding exponent \(2\) is the
minimax-optimal \textbf{guaranteed} exponent.

\subsubsection{6.3 Proof of Corollary (Large-Deviations
Refinement)}\label{proof-of-corollary-large-deviations-refinement}

For the binary symmetric case (\(K=2\), \(\theta^* = 1/2\),
\(\mu = 1/2 - \Delta\)):

\[\begin{aligned}
\text{KL}(1/2 \parallel 1/2 - \Delta) &= \frac{1}{2}\log\frac{1/2}{1/2-\Delta} + \frac{1}{2}\log\frac{1/2}{1/2+\Delta} \\
&= \frac{1}{2}\left[-\log(1-2\Delta) - \log(1+2\Delta)\right] \\
&= -\frac{1}{2}\log(1 - 4\Delta^2) \\
&= 2\Delta^2 + \frac{4}{3}\Delta^4 + O(\Delta^6)
\end{aligned}\]

By the Cramer-Chernoff theorem, for the optimal Bayes test \(\psi^*\):

\[\begin{aligned}
\lim_{M \to \infty} -\frac{1}{M}\log P_0(\psi^*=1) &= \text{KL}(1/2 \parallel \mu) = 2\Delta^2 + O(\Delta^4) \\
\lim_{M \to \infty} -\frac{1}{M}\log P_1(\psi^*=0) &= \text{KL}(1/2 \parallel 1-\mu) = 2\Delta^2 + O(\Delta^4)
\end{aligned}\]

Since \(\psi^*\) is optimal (achieves the best error rate), the minimax
risk satisfies:

\[\lim_{M \to \infty} -\frac{1}{M}\log \inf_{\psi} \max\{P_0(\psi=1), P_1(\psi=0)\} = 2\Delta^2 + O(\Delta^4)\]

As \(\Delta \to 0\), the leading constant tends to \(2\), proving the
exponent is exactly \(2\) in the small-gap (high-dimensional) limit.
\(\square\)

\begin{center}\rule{0.5\linewidth}{0.5pt}\end{center}

\subsection{7. Discussion of Tightness and
Gaps}\label{discussion-of-tightness-and-gaps}

\subsubsection{7.1 What Has Been Proved}\label{what-has-been-proved}

\begin{enumerate}
\def\labelenumi{\arabic{enumi}.}
\item
  \textbf{Testing lower bound (Part a)}: Complete for \(K=2\) using
  Slud's inequality for Binomial tails. Extended to \(K>2\) via the
  mixture-product TV bound (Lemma 4) and the fact that \(K=2\) gives the
  smallest \(\chi^2\) divergence (hardest case).
\item
  \textbf{F1 lower bound (Part b)}: Complete using
  \(1-\text{F1} \geq 1-\text{TPR}\). The bound
  \(1-\text{F1} \geq \frac{1}{2}\exp(-2M\Delta^2)\) matches the exponent
  of Theorem 1.
\item
  \textbf{Rate optimality (Part c)}: Theorem 1 (upper bound) gives
  \(\text{F1} \geq 1 - \frac{1}{\eta}\sum_s\rho_s\exp(-2M\Delta_s^2)\).
  Theorem 4 (lower bound) gives
  \(\inf_\psi\sup_P[1-\text{F1}] \geq \frac{1}{2}\exp(-2M\Delta^2)\).
  The exponent \(2M\Delta^2\) matches, confirming rate optimality.
\item
  \textbf{Large-deviations refinement}: Complete via Cramer-Chernoff,
  showing the exact asymptotic exponent is
  \(\text{KL}(\theta^* \parallel \mu) = 2\Delta^2 + O(\Delta^4)\), with
  \(2\Delta^2\) as the worst-case (minimum) exponent.
\end{enumerate}

\subsubsection{7.2 Identified Gaps}\label{identified-gaps}

\textbf{Gap 1: Mixture-product bound tightness for moderate \(M\)
(medium).} For \(K > 2\), the convexity bound
\(\text{TV}(P_0, P_1) \leq \frac{1}{L}\sum_\ell \text{TV}(P_0, Q_\ell)\)
is potentially loose. When the mixture components are far apart, the
mixture \(P_1\) is more spread out than any single component,
potentially making the TV smaller than the bound suggests. For our
specific case (all components have the same parameter under A6), the
bound is tight (Corollary 4.1).

\emph{Status}: Tight for \(C_{\text{bal}} = 1\) (balanced errors). For
\(C_{\text{bal}} > 1\), the bound may be loose but the exponential rate
is preserved.

\textbf{Gap 2: F1 lower bound tightness (minor).} The bound
\(1-\text{F1} \geq 1-\text{TPR}\) loses information about the false
positive rate. For the specific hard construction where \(P_0\) is clean
and \(P_1\) is noise, the FPR is small by design, so the bound is tight
within a factor of 2. For other constructions, FPR could be larger,
making the bound looser.

\emph{Status}: The factor \(2\) gap between \(1-\text{F1}\) and
\(1-\text{TPR}\) is unavoidable without additional assumptions.

\textbf{Gap 3: State aggregation (minor).} The lower bound is proved for
a single state \(s\) with \(\rho_s = 1\). For multi-state settings, the
worst-case over \(\mathcal{P}_\Delta\) can place all probability mass on
the hardest state, achieving the same bound. This is a standard
reduction.

\emph{Status}: Complete -- no gap.

\textbf{Gap 4: Extension to non-symmetric experts (moderate).} The proof
assumes all experts have the same error rate \(\mu\). For heterogeneous
error rates \(\mu_1,\dots,\mu_M\), the Binomial testing framework no
longer applies; we would need a Poisson Binomial lower bound (Slud's
inequality for Poisson Binomial, or a Berry-Esseen bound).

\emph{Status}: Open question. Likely the same rate holds
(\(2M\bar{\Delta}^2\) with average gap \(\bar{\Delta}\)), but the proof
would require different technical tools.

\textbf{Gap 5: Exact constant \(c = 2\) for finite \(M\) (asymptotic
only).} The large-deviations refinement is asymptotic
(\(M \to \infty\)). For finite \(M\), the exact optimal constant may be
larger or smaller than \(2\) depending on \(\mu\) and \(\Delta\). The
Slud inequality gives the finite-\(M\) lower bound with constant \(2\),
confirming that no method can achieve exponent \(> 2\) even for finite
\(M\).

\emph{Status}: The finite-\(M\) lower bound (Slud) matches the
asymptotic rate, so the exponent \(2\) holds for all \(M\). This gap is
closed.

\subsubsection{7.3 Summary}\label{summary}

\begin{longtable}[]{@{}
  >{\raggedright\arraybackslash}p{(\linewidth - 4\tabcolsep) * \real{0.3929}}
  >{\raggedright\arraybackslash}p{(\linewidth - 4\tabcolsep) * \real{0.2857}}
  >{\raggedright\arraybackslash}p{(\linewidth - 4\tabcolsep) * \real{0.3214}}@{}}
\toprule\noalign{}
\begin{minipage}[b]{\linewidth}\raggedright
Component
\end{minipage} & \begin{minipage}[b]{\linewidth}\raggedright
Status
\end{minipage} & \begin{minipage}[b]{\linewidth}\raggedright
Details
\end{minipage} \\
\midrule\noalign{}
\endhead
\bottomrule\noalign{}
\endlastfoot
\(K=2\) testing lower bound & \textbf{Complete} & Slud's inequality,
optimal Bayes test \\
\(K>2\) testing lower bound & \textbf{Complete} & Lemma 4, \(K=2\) is
hardest case \\
F1 extension & \textbf{Complete} & \(1-\text{F1} \geq 1-\text{TPR}\)
bound \\
Rate optimality & \textbf{Complete} & Exponent \(2M\Delta^2\) matches
Thm 1 \\
Large-deviations refinement & \textbf{Complete} & Cramer-Chernoff,
\(\text{KL} \to 2\Delta^2\) \\
Non-symmetric experts & \textbf{Open} & Requires Poisson Binomial
tools \\
Mixture bound for \(C_{\text{bal}}>1\) & \textbf{Partial} & Rate holds,
constant may differ \\
\end{longtable}

\subsubsection{7.4 Connection to Theorem 3
(Unidentifiability)}\label{connection-to-theorem-3-unidentifiability}

Theorem 3 shows that without assumptions (A1)-(A6), noise and difficulty
are observationally indistinguishable. Theorem 4 shows that even WITH
these assumptions, there is a fundamental statistical limit: no detector
can surpass \(\exp(-2M\Delta^2)\).

Together with Theorem 3, this forms a complete theoretical picture:

\begin{itemize}
\tightlist
\item
  \textbf{Theorem 3 (lower bound on assumptions)}: Without the SCX
  assumptions, detection is impossible (identifiability failure).
\item
  \textbf{Theorem 4 (lower bound on performance)}: With the SCX
  assumptions, detection is possible but limited by an exponential lower
  bound matching Theorem 1.
\item
  \textbf{Theorem 1 (upper bound on performance)}: The SCX detector
  achieves this optimal rate.
\end{itemize}

This completes the three-part optimality characterization: - What is
lost if assumptions fail: \textbf{all} detection ability (Theorem 3) -
What even the best detector achieves: \(\sim\exp(-2M\Delta^2)\) (Theorem
4) - What SCX actually achieves: \(\sim\exp(-2M\Delta^2)\) (Theorem 1)

\begin{center}\rule{0.5\linewidth}{0.5pt}\end{center}

\subsection{8. References}\label{references}

\begin{enumerate}
\def\labelenumi{\arabic{enumi}.}
\item
  Slud, E. V. (1977). Distribution inequalities for the binomial law.
  \emph{The Annals of Probability}, 5(3), 404-412.
\item
  Tsybakov, A. B. (2009). \emph{Introduction to Nonparametric
  Estimation}. Springer. (Section 2.1-2.4: Le Cam's method, Assouad's
  lemma, Fano's inequality)
\item
  Le Cam, L. (1973). Convergence of estimates under dimensionality
  restrictions. \emph{The Annals of Statistics}, 1(1), 38-53.
\item
  Chernoff, H. (1952). A measure of asymptotic efficiency for tests of a
  hypothesis based on the sum of observations. \emph{The Annals of
  Mathematical Statistics}, 23(4), 493-507.
\item
  Cramer, H. (1938). Sur un nouveau theoreme-limite de la theorie des
  probabilites. \emph{Actualites Scientifiques et Industrielles}, 736,
  5-23.
\item
  Hoeffding, W. (1963). Probability inequalities for sums of bounded
  random variables. \emph{Journal of the American Statistical
  Association}, 58(301), 13-30.
\item
  Yu, B. (1997). Assouad, Fano, and Le Cam. In \emph{Festschrift for
  Lucien Le Cam}, pp.~423-435. Springer.
\item
  Polyanskiy, Y., \& Wu, Y. (2022+). \emph{Information Theory: From
  Coding to Learning}. Cambridge University Press. (Forthcoming.)
\item
  Wainwright, M. J. (2019). \emph{High-Dimensional Statistics: A
  Non-Asymptotic Viewpoint}. Cambridge University Press. (Chapter 15:
  Minimax lower bounds.)
\item
  SCX Theorem 1: Multi-Expert Consistency Guarantees for Label Noise
  Detection. \texttt{../theorems/01\_noise\_detection\_guarantee.md}.
\item
  SCX Theorem 3: Unidentifiability of Noise vs.~Learnable Difficulty.
  \texttt{../theorems/03\_unidentifiability\_theorem.md}.
\item
  Minimax Optimality Feasibility Analysis.
  \texttt{./minimax\_optimality.md}.
\end{enumerate}

\begin{center}\rule{0.5\linewidth}{0.5pt}\end{center}

\subsection{Appendix A: Slud's Inequality -- Statement and Proof
Sketch}\label{appendix-a-sluds-inequality-statement-and-proof-sketch}

\subsubsection{Statement}\label{statement}

\textbf{Theorem (Slud, 1977, Theorem 2).} Let
\(S_M \sim \text{Binomial}(M, p)\) with \(p \leq 1/2\). Then for any
integer \(k\) with \(Mp \leq k \leq M/2\):

\[\mathbb{P}(S_M \geq k) \geq \frac{1}{2} \cdot \exp\!\left(-2M\left(\frac{k}{M} - p\right)^2\right)\]

This is a lower bound (reverse inequality) for the upper tail of the
Binomial distribution. The factor \(1/2\) is universal (not depending on
\(M, p, k\)), and the exponential rate \(2M(k/M - p)^2\) matches the
Hoeffding upper bound.

\subsubsection{Proof Sketch}\label{proof-sketch}

Slud's proof uses the representation of the Binomial CDF as a
regularized incomplete beta function:

\[\mathbb{P}(S_M \geq k) = I_p(k, M-k+1) = \frac{B(p; k, M-k+1)}{B(k, M-k+1)}\]

where \(B(p; a, b)\) is the incomplete beta function. The key inequality
follows from the monotonicity of the likelihood ratio for Binomial
distributions (Karlin's totally positive property) and an integral
bound.

An alternative modern proof uses the following martingale argument (due
to Shorack \& Wellner, 1986, pp.~440-441): For \(p \leq 1/2\), the
Binomial probability can be bounded below by the Gaussian approximation
times a factor that depends on \(p(1-p)\). The worst case \(p = 1/2\)
gives the factor \(1/2\).

\subsubsection{Why the Factor 1/2 is
Important}\label{why-the-factor-12-is-important}

The factor \(1/2\) in Slud's inequality is what makes the lower bound
non-trivial. Without it, the bound could be arbitrarily close to \(0\)
(which is always true for any probability). With it, we guarantee that
the Binomial tail probability is at least half the Hoeffding bound,
establishing that the Hoeffding exponent is tight (up to constant
factors).

\subsubsection{Connection to Central Limit
Theorem}\label{connection-to-central-limit-theorem}

For large \(M\), the de Moivre--Laplace theorem gives:

\[\mathbb{P}(S_M \geq M/2) \approx 1 - \Phi\!\left(\frac{(M/2) - M\mu}{\sqrt{M\mu(1-\mu)}}\right) = 1 - \Phi\!\left(-\frac{2\Delta\sqrt{M}}{\sqrt{4\mu(1-\mu)}}\right)\]

where \(\Phi\) is the standard normal CDF. Using the Mills ratio bound
\(1 - \Phi(-x) \geq \frac{x}{1+x^2} \cdot \frac{e^{-x^2/2}}{\sqrt{2\pi}}\)
for \(x > 0\), the exponent is \(x^2/2 = 2M\Delta^2/(4\mu(1-\mu))\),
which is \(\geq 2M\Delta^2\) since \(4\mu(1-\mu) \leq 1\). This provides
an alternative (asymptotic) verification of the \(2M\Delta^2\) exponent.

\begin{center}\rule{0.5\linewidth}{0.5pt}\end{center}

\subsection{Appendix B: Numerical
Illustration}\label{appendix-b-numerical-illustration}

For concrete parameter values, the lower bound on testing error gives:

\begin{longtable}[]{@{}cccc@{}}
\toprule\noalign{}
\(\Delta\) & \(\mu\) & \(M\) & Lower bound
\(\frac{1}{2}\exp(-2M\Delta^2)\) \\
\midrule\noalign{}
\endhead
\bottomrule\noalign{}
\endlastfoot
0.3 & 0.2 & 10 & \(0.5 \cdot e^{-1.8} \approx 0.083\) \\
0.3 & 0.2 & 20 & \(0.5 \cdot e^{-3.6} \approx 0.014\) \\
0.25 & 0.25 & 30 & \(0.5 \cdot e^{-3.75} \approx 0.012\) \\
0.2 & 0.3 & 50 & \(0.5 \cdot e^{-4.0} \approx 0.009\) \\
0.1 & 0.4 & 100 & \(0.5 \cdot e^{-2.0} \approx 0.068\) \\
\end{longtable}

Note that the lower bound is non-trivial (below \(0.1\)) when
\(M\Delta^2\) is above approximately \(1\). The condition
\(M\Delta^2 > 1\) corresponds to the regime where detection is possible
with non-trivial accuracy.

For comparison, the Theorem 1 upper bound on \(1-\text{F1}\) (assuming
\(\eta=0.1\), \(\rho_s=1\)) gives:

\[\text{Upper bound on } 1-\text{F1} \leq \frac{1}{\eta}\exp(-2M\Delta^2) = 10\exp(-2M\Delta^2)\]

The gap between upper and lower bounds is a factor of \(20/\eta\) in the
pre-exponential constant, but the exponent \(2M\Delta^2\) is identical.

\begin{center}\rule{0.5\linewidth}{0.5pt}\end{center}

\subsection{Appendix C: Adversarial Construction for the Lower
Bound}\label{appendix-c-adversarial-construction-for-the-lower-bound}

The proof of Theorem 4 constructs a pair \((P_0, P_1)\) within the class
\(\mathcal{P}_\Delta\). Here we verify that this construction satisfies
all SCX assumptions:

\begin{longtable}[]{@{}
  >{\raggedright\arraybackslash}p{(\linewidth - 2\tabcolsep) * \real{0.4800}}
  >{\raggedright\arraybackslash}p{(\linewidth - 2\tabcolsep) * \real{0.5200}}@{}}
\toprule\noalign{}
\begin{minipage}[b]{\linewidth}\raggedright
Assumption
\end{minipage} & \begin{minipage}[b]{\linewidth}\raggedright
Verification
\end{minipage} \\
\midrule\noalign{}
\endhead
\bottomrule\noalign{}
\endlastfoot
\textbf{A1 (Disjoint training)} & Experts trained on disjoint subsets of
clean data \\
\textbf{A2 (Cond. independence)} & Given \(x\), \(\{e_m\}\) are i.i.d.
Bernoulli \\
\textbf{A3 (Bounded loss)} & 0-1 loss with \(B=1\) \\
\textbf{A4 (Uniform noise)} & Noise label uniform over
\(\mathcal{Y}\setminus\{y^*\}\) \\
\textbf{A5 (State homogeneity)} & Single state \(s\) with constant
\(\mu\) \\
\textbf{A6 (Balanced errors)} & \(C_{\text{bal}} = 1\) (errors uniform
across wrong classes) \\
\end{longtable}

The construction uses the most favorable settings within the assumption
class, making the detection problem as hard as possible. This is the
standard approach for minimax lower bounds: find the hardest problem
instance within the class.

Since
\(\inf_\psi \sup_P \geq \inf_\psi \max\{P_0(\psi=1), P_1(\psi=0)\}\)
(the supremum over \(P\) is at least the value at any specific \(P\)),
the lower bound holds.

\begin{center}\rule{0.5\linewidth}{0.5pt}\end{center}

\textbf{End of proof document.}

\end{document}
